% =============================================================
%  Chapter 2 -- Thermal-history modeling of Win-06-01b
%  (methods + results). Body text only: \input from a wrapper.
%  Requires: amsmath, graphicx, natbib-style \citep/\citet,
%  % =============================================================
%  chapter2_derived_values.tex -- AUTO-GENERATED, DO NOT EDIT BY HAND.
%  Written by Figures/write_chapter2_macros.py from
%  Figures/generated/data/stage1_regression_cache.json (itself written by
%  Figures/prepare_regression_cache.py from the master workbook).
%
%  Regenerate after any data/regression change:
%    python3 Figures/prepare_regression_cache.py
%    python3 Figures/write_chapter2_macros.py
%  (run_all_figures.py does both automatically.)
%
%  \input{} by BOTH chapter2goethite.tex (standalone) and
%  Chapter2/chapter2.tex (dissertation chapter) so the Abstract,
%  Results, Discussion, and Conclusions all read the same numbers.
%
%  Stage 1 reference line (3He, "treatment as uncertainty" central/M
%  bound, fractional 3He yield >= 0.07, n=12 points, R2=0.9890):
%    ChTwoEa/ChTwoEaErr : activation energy (kJ/mol)
%    ChTwoTc/ChTwoTcErr : Dodson (1973) closure temperature (degC),
%      intercept-only regression uncertainty (D0/a^2 term only, Ea
%      held fixed) -- see thermochr_utils.dodson_closure_temperature_stderr_interceptonly.
%      Joint slope+intercept uncertainty is +/- 14.9 degC if ever needed.
%    ChTwoRtwo/ChTwoN   : regression R^2 and number of points
%
%  Isothermal 3He permeation Arrhenius regression (from
%  generated/data/naive_lambda_arrhenius_cache.json, n=10 sequences, R2=0.9362):
%    ChTwoPermEa/ChTwoPermEaErr : activation energy, 2 s.e. (cal/mol)
%    ChTwoPermRtwo/ChTwoPermN   : regression R^2 and number of sequences
%    ChTwoPermLamMin/ChTwoPermLamMax : range of the 10 per-series fitted
%      rate constants themselves (s^-1), math-mode, e.g. "$1.2\times10^{-6}$"
% =============================================================
\newcommand{\ChTwoEa}{157.1}
\newcommand{\ChTwoEaErr}{5.3}
\newcommand{\ChTwoTc}{98.1}
\newcommand{\ChTwoTcErr}{8.6}
\newcommand{\ChTwoRtwo}{0.989}
\newcommand{\ChTwoN}{12}
\newcommand{\ChTwoPermEa}{10192}
\newcommand{\ChTwoPermEaErr}{1882}
\newcommand{\ChTwoPermRtwo}{0.936}
\newcommand{\ChTwoPermN}{10}
\newcommand{\ChTwoPermLamMin}{$2.7\times10^{-5}$}
\newcommand{\ChTwoPermLamMax}{$6.3\times10^{-5}$}
 and
%  % Auto-generated by Chapter2/Figures/fig_winsor_burial_path.py -- do not edit by hand.
\newcommand{\ChTwoStageOneLnDzeroAtwo}{16.16}
\newcommand{\ChTwoLRDfrac}{1.7}
\newcommand{\ChTwoLRDsize}{0.0024}
\newcommand{\ChTwoLRDmisfit}{7}
\newcommand{\ChTwoSingleDomainMisfit}{44}
\newcommand{\ChTwoBurialNfit}{24}
\newcommand{\ChTwoBurialTmaxLo}{60}
\newcommand{\ChTwoBurialTmaxHi}{64}
\newcommand{\ChTwoBurialTpeakLo}{57.5}
\newcommand{\ChTwoBurialTpeakHi}{87.5}
\newcommand{\ChTwoBurialTformLo}{180}
\newcommand{\ChTwoBurialTformHi}{200}
\newcommand{\ChTwoBurialBestTmax}{62}
\newcommand{\ChTwoBurialBestTpeak}{77.5}
\newcommand{\ChTwoBurialBestTform}{200}
\newcommand{\ChTwoBurialBestAge}{179}
\newcommand{\ChTwoBurialBestChi}{2.3}
\newcommand{\ChTwoBurialGlobalChi}{2.2}
\newcommand{\ChTwoBurialChiTformLo}{2.18}
\newcommand{\ChTwoBurialChiTformHi}{2.42}
\newcommand{\ChTwoVTmax}{60}
\newcommand{\ChTwoVChi}{2.5}
\newcommand{\ChTwoVAge}{180}
\newcommand{\ChTwoIsoChi}{16}
 in the preamble
%  (the latter is written by Figures/fig_winsor_burial_path.py).
%  Figures: Figures/fig_winsor_3he_lrd_fit.py and
%           Figures/fig_winsor_burial_path.py.
% =============================================================

\subsection{Thermal-History Modeling of Win-06-01b}\label{sec:gt_thermal_methods}

Win-06-01b was collected as float in the Flinders Ranges and has a
(U--Th)/He age of 130--180~Ma \citep{waltenberg2012}. Apatite fission-track
studies across the region record regional cooling from roughly 50--70\,\textdegree C
after the mid-Cretaceous \citep{gibsonstuwe2000,mitchell1998}, and Triassic
intramontane basins less than 1.5~km thick preserve near-surface Early Mesozoic
sections \citep{mitchell1998}. We therefore ask which burial-and-exhumation
histories reproduce both the bulk-normalized $^4$He/$^3$He spectrum of the
water-saturated experiment (Fig.~\ref{fig:ch2-ratio-spectra}) and the He age,
following the ratio-evolution approach of \citet{shusterfarley2004}.

\paragraph{Diffusion domains.}
Helium resides in two spherical domains of the same goethite, with a common
activation energy and frequency factor: a high-retentivity domain (HRD) of
radius $a$ and a low-retentivity domain (LRD) of radius $r a$ ($r<1$) that holds
a fraction $f$ of the uniformly distributed, proton-induced $^3$He. The HRD
follows the Stage~1 Arrhenius regression of the water-saturated experiment
(Fig.~\ref{fig:ch2-arrhenius-zoom}; $E_a = \ChTwoEa$~kJ\,mol$^{-1}$,
$\ln(D_0/a^2) = \ChTwoStageOneLnDzeroAtwo$, $D_0/a^2$ in s$^{-1}$), so that
\begin{equation}
  \frac{D}{a^2}(T) = \frac{D_0}{a^2}\exp\!\left(-\frac{E_a}{RT}\right),
  \qquad
  \left(\frac{D}{a^2}\right)_{\mathrm{LRD}} = r^{-2}\,\frac{D}{a^2}.
  \label{eq:gt_th_arrhenius}
\end{equation}
Diffusion in each domain is described by the non-dimensional time
\begin{equation}
  \tau(t) = \int \frac{D\big(T(t')\big)}{a^2}\,\mathrm{d}t',
  \label{eq:gt_th_tau}
\end{equation}
and an initially uniform sphere releases the cumulative fraction
\citep{crank1979}
\begin{equation}
  F(\tau) = 1 - \frac{6}{\pi^2}\sum_{k=1}^{\infty}\frac{1}{k^2}\exp\!\left(-k^2\pi^2\tau\right).
  \label{eq:gt_th_F}
\end{equation}

\paragraph{Constraining the LRD with $^3$He.}
With the HRD's laboratory time $\tau_i$ after step $i$ accumulated from
Eq.~\ref{eq:gt_th_arrhenius} and the heating schedule, the model cumulative
$^3$He release is
\begin{equation}
  \Sigma F^{\mathrm{model}}_{3,i} = f\,F\!\left(\tau_i/r^2\right) + (1-f)\,F(\tau_i).
  \label{eq:gt_th_twodomain}
\end{equation}
We fit $f$ and $r$ by minimizing
$\sum_i \big[(\Sigma F^{\mathrm{model}}_{3,i}-\Sigma F^{\mathrm{meas}}_{3,i})/\Sigma F^{\mathrm{meas}}_{3,i}\big]^2$
over all 34 steps, with $f \le 0.075$ as required by the delta plot
(Fig.~\ref{fig:ch2-delta-zoom}). These two parameters are then held fixed, so
the $^4$He/$^3$He spectrum constrains only the thermal history.

\paragraph{Radiogenic $^4$He in nature.}
In each domain, $^4$He is produced uniformly at rate $P$ and diffuses with a
zero-concentration boundary,
\begin{equation}
  \frac{\partial C}{\partial t} = \frac{D}{a^2}\,\frac{1}{x^2}\frac{\partial}{\partial x}\!\left(x^2\frac{\partial C}{\partial x}\right) + P,
  \qquad C(1,t)=0,\quad C(x,t_\mathrm{form})=0,
  \label{eq:gt_th_pde}
\end{equation}
where $x$ is radial position normalized to the domain radius and
$t_\mathrm{form}$ is the crystallization age. Expanding the concentration in
eigenfunctions \citep{carslawjaeger1959},
$C(x,t)=\sum_k A_k(t)\,\sin(k\pi x)/x$, each mode obeys
\begin{equation}
  \frac{\mathrm{d}A_k}{\mathrm{d}t} = -k^2\pi^2\frac{D}{a^2}A_k + c_k P,
  \qquad c_k = \frac{2(-1)^{k+1}}{k\pi},
  \label{eq:gt_th_modes}
\end{equation}
where $c_k$ are the coefficients of a uniform profile. Discretizing the
thermal history into intervals $\Delta t$ (2~Myr) of constant temperature gives
the exact update
\begin{equation}
  A_k \leftarrow A_k\,e^{-\lambda_k} + c_k P\,\Delta t\,\frac{1-e^{-\lambda_k}}{\lambda_k},
  \qquad \lambda_k = k^2\pi^2\,\frac{D}{a^2}\,\Delta t ,
  \label{eq:gt_th_update}
\end{equation}
which we evaluate with 8000 modes. The $^4$He content per unit volume is
$N = 3\sum_k A_k\,(-1)^{k+1}/(k\pi)$, and, taking U and Th to be distributed
like the $^3$He, the predicted bulk He age is
\begin{equation}
  t_\mathrm{He} = \frac{(1-f)\,N_\mathrm{HRD} + f\,N_\mathrm{LRD}}{P},
  \label{eq:gt_th_age}
\end{equation}
neglecting parent decay (a $\sim$3\% effect over 200~Myr).

\paragraph{Laboratory degassing of the natural profile.}
In the laboratory each mode decays as $\exp(-k^2\pi^2\tau)$. The HRD's
laboratory times $\tau_i$ are solved from the measured cumulative $^3$He
through Eq.~\ref{eq:gt_th_twodomain} (the LRD's are $\tau_i/r^2$), which removes
any dependence on the heating schedule or on step-to-step kinetic scatter. The
predicted spectrum is
\begin{equation}
  \frac{R_{\mathrm{step},i}}{R_\mathrm{bulk}} =
  \frac{\Delta N_{4,i}/N_{4,\mathrm{tot}}}{\Delta F_{3,i}},
  \label{eq:gt_th_ratio}
\end{equation}
where $\Delta N_{4,i}$ is the $^4$He released from both domains in step $i$
and $\Delta F_{3,i}$ is the measured fraction of total $^3$He released in that
step.

\paragraph{Thermal histories and fitting.}
The thermal history is piecewise linear through
$(t_\mathrm{form}, T_s)$, $(t_\mathrm{peak}, T_\mathrm{max})$, $(50~\mathrm{Ma}, T_s)$,
and $(0, T_s)$, with $T_s = 20$\,\textdegree C: the goethite crystallizes at
the surface, is heated during burial to a peak temperature, and is exhumed by
the Paleogene, when Paleogene sediments in Flinders Ranges valleys show that
basement was exposed \citep{mitchell1998}. We searched $t_\mathrm{form}$ from
180 to 210~Ma (5~Myr steps), $t_\mathrm{peak}$ from 55~Ma to $t_\mathrm{form}$
(2.5~Myr steps), and $T_\mathrm{max}$ from 20 to 90\,\textdegree C (1\,\textdegree C
steps). Each history is scored against the $\ChTwoBurialNfit$ steps with
lag-correction bounds and $>0.1$~pccSTP $^3$He by
\begin{equation}
  \chi^2 = \sum_i\left(\frac{R^{\mathrm{model}}_i - R^{\mathrm{M}}_i}{\sigma_i}\right)^2,
  \qquad \sigma_i = \max\!\left[\tfrac{1}{2}\left(R^{\mathrm{UB}}_i-R^{\mathrm{LB}}_i\right),\ 0.02\right],
  \label{eq:gt_th_chi2}
\end{equation}
where M, LB and UB are the central, lower-bound and upper-bound
bulk-normalized ratios of the bounding method. Histories within
$\Delta\chi^2 \le 3.53\,\max(1,\chi^2_\nu)$ of the minimum form the joint
1$\sigma$ region for three parameters. For comparison we also fit isothermal
holding (one parameter, $\tau_\mathrm{nat}$) and a symmetric history heating
linearly from 200 to 100~Ma and cooling linearly to the present (one
parameter, $T_\mathrm{max}$). The kinetic extrapolation uncertainty of the
Stage~1 line ($\pm$$\sim$4\,\textdegree C near 60\,\textdegree C) is not
propagated.

% -------------------------------------------------------------
\subsection{Thermal History of Win-06-01b}\label{sec:gt_thermal_results}

The two-domain model reproduces the measured cumulative $^3$He release to
$\ChTwoLRDmisfit$\% (root-mean-square relative misfit), compared with
$\ChTwoSingleDomainMisfit$\% for the Stage~1 line alone. It requires an LRD
holding $\ChTwoLRDfrac$\% of the $^3$He with $r = \ChTwoLRDsize$, which is
degassed within the first $\sim$7 steps and holds negligible radiogenic $^4$He
in nature (Figs.~\ref{fig:gt_th_arrhenius} and \ref{fig:gt_th_delta}).

With the LRD fixed, burial histories peaking at
$\ChTwoBurialTmaxLo$--$\ChTwoBurialTmaxHi$\,\textdegree C between
$\ChTwoBurialTpeakHi$ and $\ChTwoBurialTpeakLo$~Ma fit the spectrum
($\chi^2_\nu \approx \ChTwoBurialGlobalChi$--$\ChTwoBurialBestChi$) and predict
He ages within the measured range for crystallization at
$\ChTwoBurialTformLo$--$\ChTwoBurialTformHi$~Ma
(Figs.~\ref{fig:gt_th_map}--\ref{fig:gt_th_paths}). The spectrum alone barely
distinguishes crystallization ages ($\chi^2_\nu = \ChTwoBurialChiTformLo$--$\ChTwoBurialChiTformHi$
across 180--210~Ma); the He age excludes crystallization before
$\sim$\ChTwoBurialTformHi~Ma. The best-fitting age-consistent history crystallizes
at $\ChTwoBurialBestTform$~Ma, reaches $\ChTwoBurialBestTmax$\,\textdegree C at
$\ChTwoBurialBestTpeak$~Ma and predicts a He age of $\ChTwoBurialBestAge$~Ma.
These peak temperatures and timing agree with the 50--70\,\textdegree C
cooling after $\sim$85~Ma inferred from apatite fission tracks
\citep{gibsonstuwe2000} and sit at the $<$60\,\textdegree C limit reported
near Copley \citep{mitchell1998}. The symmetric history fits almost as well
($T_\mathrm{max} = \ChTwoVTmax$\,\textdegree C, $\chi^2_\nu = \ChTwoVChi$,
$t_\mathrm{He} = \ChTwoVAge$~Ma), whereas isothermal holding is rejected
($\chi^2_\nu = \ChTwoIsoChi$; Fig.~\ref{fig:gt_th_spectrum}).

\begin{figure}[htbp]
\centering
\includegraphics[width=0.8\textwidth]{winsor_3he_lrd_arrhenius.pdf}
\caption{Arrhenius plot of $^3$He diffusivities in HT-Win-06-01b with the
two-domain model. Black circles are bounding-method central values with
lower and upper bounds; hollow circles are low-signal steps or steps without
lag bounds. Blue diamonds are apparent diffusivities obtained by applying the
same single-sphere inversion \citep{fechtigkalbitzer1966} to the model's
cumulative release (Eq.~\ref{eq:gt_th_twodomain}); the dashed line is the
Stage~1 regression, which describes the HRD. An LRD holding $\ChTwoLRDfrac$\%
of the $^3$He with $a_\mathrm{LRD}/a_\mathrm{HRD} = \ChTwoLRDsize$ reproduces
the elevated early diffusivities and the retrograde steps.}
\label{fig:gt_th_arrhenius}
\end{figure}

\begin{figure}[htbp]
\centering
\includegraphics[width=0.8\textwidth]{winsor_3he_lrd_delta.pdf}
\caption{Delta plot of the same data and model: departure of apparent
$\ln(D/a^2)$ from the Stage~1 line against cumulative $^3$He fraction. The
model approaches the Stage~1 line as the LRD is exhausted, consistent with
an LRD holding no more than 7.5\% of the $^3$He.}
\label{fig:gt_th_delta}
\end{figure}

\begin{figure}[htbp]
\centering
\includegraphics[width=\textwidth]{winsor_burial_chi2_map.pdf}
\caption{Fit of burial-and-exhumation histories to the $^4$He/$^3$He
spectrum for crystallization at 180, 195 and 210~Ma. Colors show reduced
$\chi^2$ as a function of the time and magnitude of peak temperature; white
outlines enclose the joint 1$\sigma$ region across all three parameters, and
diamonds mark each panel's best fit. Dashed lines are isopleths of the
predicted bulk He age at the limits of the measured range (130 and 180~Ma).
Orange boxes show the peak window implied by apatite fission-track data
\citep{gibsonstuwe2000,mitchell1998}.}
\label{fig:gt_th_map}
\end{figure}

\begin{figure}[htbp]
\centering
\includegraphics[width=0.8\textwidth]{winsor_burial_spectrum.pdf}
\caption{Bulk-normalized $^4$He/$^3$He spectrum of HT-Win-06-01b (black,
with bounding-method LB--UB boxes) compared with the best-fitting
age-consistent burial history (blue), the best symmetric 200--100--0~Ma
history (orange) and the best isothermal history (green), all computed with
the LRD fixed from the $^3$He fit.}
\label{fig:gt_th_spectrum}
\end{figure}

\begin{figure}[htbp]
\centering
\includegraphics[width=0.8\textwidth]{winsor_burial_tT_paths.pdf}
\caption{Thermal histories within the joint 1$\sigma$ region of the
spectrum fit. Blue paths also predict He ages within the measured
130--180~Ma range; gray paths predict older ages. The black path is the
best-fitting age-consistent history, and the orange box is the apatite
fission-track peak window.}
\label{fig:gt_th_paths}
\end{figure}
